%% kapa: exam=1 round=26 year=2015 period=1 form=A subject=civil_law subject_ko=민법 seq=11 number=11
%% source: Q-Net 감정평가사 기출문제 (cst00309) artlSeq=5138001 file=extracted/1008/5138001/제26회_감정평가사 1차_기출/제26회_감정평가사 1차_기출/1교시_감평_A형.pdf
%% conversion: kordoc 4.24.51 -> tools/gen1.py (auto); review=proofread-sonnet; flags=none
\begin{question}{11}
무권대리에 관한 설명으로 옳지 않은 것은? (다툼이 있으면 판례에 따름)
\begin{choices}
  \item 무권대리행위에 대한 본인의 추인은 재판상ㆍ재판 외에서 묵시적으로도 할 수 있다.
  \item 무권대리행위는 추인이나 거절 전에는 유동적 무효이다.
  \item 본인이 무권대리행위의 추인을 거절한 후에는 다시 추인할 수 없다.
  \item 매매계약을 체결한 무권대리인으로부터 매매대금의 일부를 본인이 수령한 경우 특별한 사정이 없는 한 본인이 무권대리행위를 묵시적으로 추인한 것으로 본다.
  \item 무권대리인의 상대방에 대한 책임은 과실책임이다.
\end{choices}
\end{question}
\answer{5}
